\chapter{Transactions, isolation, and concurrent work}
\label{ops:transactions}

\section{What a transaction gives an application}

A transaction groups changes into one unit that can be committed or rolled
back. Consider an order: its header, line items, and stock adjustment should
either become durable together or leave the database unchanged. A sequence
of successful individual statements is not enough if the application fails
between them.

MiniSQL uses autocommit when no explicit transaction is active. Each ordinary
data-changing statement receives its own transaction. To group statements,
use \code{BEGIN}, execute the related operations, and finish with
\code{COMMIT}. Use \code{ROLLBACK} when the operation cannot be completed.

\begin{lstlisting}[language=SQL]
BEGIN;
INSERT INTO orders(order_id, customer_id, order_status)
VALUES (900, 1, 'new');
INSERT INTO order_items(order_id, product_id, quantity, unit_price)
VALUES (900, 101, 2, 79.90);
UPDATE products SET stock = stock - 2 WHERE product_id = 101;
COMMIT;
\end{lstlisting}

This example illustrates transaction boundaries; an application must still
ensure that the product exists and enough stock is available. A transaction
does not invent a business rule. Express an invariant through a constraint
where possible, or check the affected-row count of a conditional update
inside the transaction.

\begin{lstlisting}[language=SQL]
UPDATE products
SET stock = stock - 2
WHERE product_id = 101 AND stock >= 2;
\end{lstlisting}

If the affected-row count is zero, the application should not proceed as
though stock were reserved. It should roll back the whole order operation
or recover to an earlier savepoint according to its intended workflow.

\sourcecode{src/minisql/executor/executor.ml}{beginExplicit, executeDml, commitExplicit, and rollbackExplicit.}

\section{Complete transaction syntax}
\index{BEGIN}\index{COMMIT}\index{ROLLBACK}

\begin{lstlisting}[language=SQL]
BEGIN [TRANSACTION]
  [READ ONLY | READ WRITE]
  [ISOLATION LEVEL SERIALIZABLE | ISOLATION LEVEL READ COMMITTED];
COMMIT [TRANSACTION];
ROLLBACK [TRANSACTION];
SAVEPOINT savepoint_name;
ROLLBACK [TRANSACTION] TO [SAVEPOINT] savepoint_name;
RELEASE [SAVEPOINT] savepoint_name;
\end{lstlisting}

Square brackets in this reference denote optional syntax; do not type the
brackets. The default is \code{READ WRITE} with \code{SERIALIZABLE}
isolation. The parser accepts the read mode and isolation clause in either
order. Use each at most once in clear application SQL. There is no nested
\code{BEGIN}, no SQL \code{SET TRANSACTION}, and no
\code{REPEATABLE READ} or \code{READ UNCOMMITTED} isolation level.

\code{COMMIT} or \code{ROLLBACK} without an explicit transaction reports a
transaction-state error. Starting another explicit transaction while one is
already active also reports an error. Transaction state belongs to the
session, so sending \code{BEGIN} on one connection and the changes on
another cannot group those changes.

\sourcecode{src/minisql/sql/parser.ml}{parseBegin and transaction statement dispatch.}

\section{Read-only work and the two isolation levels}
\index{isolation}\index{READ COMMITTED}\index{SERIALIZABLE}\index{read-only transactions}

\begin{lstlisting}[language=SQL]
BEGIN READ ONLY ISOLATION LEVEL SERIALIZABLE;
SELECT COUNT(*) AS customer_count FROM customers;
SELECT COUNT(*) AS order_count FROM orders;
COMMIT;
\end{lstlisting}

A read-only transaction declares that it will not modify data. MiniSQL opens
the page transaction lazily; a SELECT-only explicit transaction need not
allocate a WAL transaction or object identifier. This is useful for reporting,
but it is not a declaration that the report will never delay a writer.

\warning{NEXTVAL is outside the ordinary read-only guarantee}{NEXTVAL advances
persisted sequence state directly, outside the page transaction's read-only
check, and rollback does not undo consumption. The native validation for this
manual confirmed a successful NEXTVAL inside BEGIN READ ONLY and a surviving
CURRVAL after ROLLBACK. The streaming cursor probe materializes expressions
before deciding whether a plan is streamable; a fallback can materialize
NEXTVAL again. Plain EXPLAIN also materializes it. Keep NEXTVAL out of
read-only reports, standby queries, and diagnostic EXPLAIN statements in this
revision. Do not assume that a SELECT-looking statement is free of side
effects or consumes one sequence value per output row. The SQL reference
documents the concrete sequence limitations.}

\begin{longtable}{p{.25\linewidth}p{.68\linewidth}}
\toprule
Isolation & Current locking behavior and application consequence \\
\midrule
\endhead
\code{READ COMMITTED} & A statement acquires a read lease and releases a
newly acquired read lease when that statement finishes. Separate SELECTs can
therefore observe changes committed between them. A read lease already held
by the transaction, or its writer ownership, is not independently released
by this statement cleanup. \\
\code{SERIALIZABLE} & Read locks are retained through the explicit transaction.
Other transactions cannot freely interleave writes that would invalidate the
protected read work. Keep the transaction short because its read ownership
can delay a writer. \\
\bottomrule
\end{longtable}

These are lock-based semantics. MiniSQL does not implement a multiversion
snapshot engine in which long readers can assume that writers continue
unimpeded. A practical diagnostic experiment uses two sessions: let session A
begin a serializable transaction and read a table; let session B attempt an
update; then commit A. B may have to wait, or reach its deadline if A remains
open too long. Repeat with read committed and observe that the read lease is
released after the statement. Use a disposable database and avoid interpreting
wall-clock timing as a correctness proof.

\sourcecode{src/minisql/transaction/lock_manager.ml}{acquireStatementRead, finishStatement, and finishTransaction.}
\sourcecode{src/minisql/executor/executor.ml}{statementIsolation and statement read-lease handling.}

\section{Savepoints: recover part of a larger operation}
\index{SAVEPOINT}\index{RELEASE}

A savepoint records a recoverable position inside an explicit DML
transaction. It is not an independent commit. Changes before and after a
savepoint remain part of the same enclosing transaction until the final
\code{COMMIT}.

\begin{lstlisting}[language=SQL]
BEGIN;
UPDATE customers SET city = 'Hamburg' WHERE customer_id = 1;
SAVEPOINT optional_change;
UPDATE products SET stock = stock + 5 WHERE product_id = 101;
ROLLBACK TO SAVEPOINT optional_change;
RELEASE SAVEPOINT optional_change;
COMMIT;
\end{lstlisting}

The city change is committed; the stock increase is undone. Rolling back to
a savepoint retains the named savepoint but removes savepoints created after
it. Releasing a savepoint removes that savepoint and later ones, without
undoing their data changes. If a name is reused, lookup selects the most
recent matching savepoint. Unique descriptive names make this stack behavior
easier to reason about.

Creating a savepoint requires an explicit transaction. It selects the DML
transaction path even if no data statement has yet run. DDL savepoints are
not supported. A rollback to a valid DML savepoint also clears the explicit
transaction's failed state, allowing a controlled recovery after a statement
error.

\subsection{Recovering from a constraint error}

Suppose a business operation has an optional insert that might collide with
an existing key. In an interactive session:

\begin{lstlisting}[language=SQL]
BEGIN;
SAVEPOINT before_optional;
INSERT INTO customers(customer_id, name)
VALUES (1, 'Duplicate identifier');
-- If the insert fails because customer 1 exists:
ROLLBACK TO before_optional;
-- Continue only after inspecting that error.
UPDATE customers SET active = TRUE WHERE customer_id = 1;
COMMIT;
\end{lstlisting}

This is a branching application workflow, not a promise that a script runner
will continue after the error. A script client normally stops on an error;
application code should catch it and issue the recovery command deliberately.
For routine insert conflicts, the SQL \code{ON CONFLICT} forms may be a
clearer design than deliberately triggering an error.

\sourcecode{src/minisql/transaction/transaction.ml}{Savepoint stack and change snapshots.}
\sourcecode{src/minisql/executor/executor.ml}{executeSavepoint, executeRollbackTo, and executeReleaseSavepoint.}

\section{The failed-transaction state}
\index{failed transactions}

When a statement error marks an explicit transaction failed, subsequent
normal work and \code{COMMIT} cannot be treated as successful continuation.
The diagnostic explicitly requests rollback. Use \code{ROLLBACK} to abandon
the transaction, or \code{ROLLBACK TO} a valid DML savepoint when the program
has a defined partial-recovery strategy. Merely retrying the failed SQL
without repairing transaction state is insufficient.

A disconnected session is cleaned up by the engine, which abandons its
uncommitted work and releases locks. Nevertheless, a lost connection during
commit creates an application-level uncertainty: the durable commit may have
completed even though its acknowledgement was not received. Reconnecting and
blindly repeating a non-idempotent operation can duplicate effects. Use a
stable application request identifier or inspect the intended business record
before retrying an operation whose commit outcome is unknown.

Distinguish three cases in application error handling: a rejected statement
whose transaction can be rolled back, a transient lock/deadline conflict
whose entire business transaction can be retried, and an uncertain
post-disconnection commit whose outcome must be reconciled. The server's
transaction machinery solves durable grouping; it cannot decide which
business request a replacement connection is trying to repeat.

\section{DDL is not uniformly transactional}

MiniSQL has a schema-journal path for a subset of table and index DDL. That
does not mean that arbitrary schema changes can be mixed with data updates
inside one transaction.

\begin{longtable}{p{.34\linewidth}p{.59\linewidth}}
\toprule
Operation & Explicit-transaction rule \\
\midrule
\endhead
Ordinary INSERT, UPDATE, DELETE, MERGE, TRUNCATE & Supported on the DML path.
Constraints, indexes, and supported triggers participate in the operation. \\
Table/index DDL handled by the schema journal & The supported staged DDL path
is available only to trusted sessions. It must not be mixed with DML. \\
SELECT after staged DDL & Rejected: the executor does not query a partly
staged schema. Commit the DDL first. \\
Authenticated table/index DDL & Autocommit only. \\
Schema, view, trigger, sequence, and procedure DDL & Autocommit only. \\
Users, roles, GRANT, and REVOKE & Autocommit only. \\
ANALYZE, VACUUM, REINDEX & Autocommit maintenance operations. \\
Savepoints for schema changes & Unsupported. \\
\bottomrule
\end{longtable}

For application migrations, the clearest workflow is therefore to execute
schema changes as explicit migration steps outside a DML transaction, verify
the new schema, and then run data transformations in their own transactions.
Keep a migration log in the application if multiple steps must be tracked.
Do not copy a PostgreSQL or MySQL migration file under the assumption that all
DDL has the same transactional behavior.

\sourcecode{src/minisql/executor/executor.ml}{ensureExplicitDml, ensureExplicitDdl, executeDdl, and object-specific DDL handlers.}

\section{How concurrent sessions are scheduled}

The server has a dedicated acceptor and a bounded native thread pool. A
worker owns one connection and its protocol/session state for the lifetime
of that connection. Independent clients can therefore parse network frames
and communicate concurrently; one slow socket does not serialize every
other socket.

Inside each open database, a writer-prioritized readers/writer execution
gate protects physical execution. Ordinary read-only SELECT, EXPLAIN, and
metadata work can share the gate. Data changes, schema changes, security
changes, maintenance, and conservative session-state
operations require exclusive execution. When a writer is waiting, the gate
stops admitting new readers until that writer gets its turn. Existing readers
finish first.

This gives real parallel execution for independent reads, while preserving
one physical writer per database. Increasing \code{maxConnections} increases
connection capacity; it does not turn the engine into a multiwriter storage
system. A workload with many idle pooled connections can occupy workers
without producing useful parallelism. Size connection pools together with
the server limit and query working memory.

The logical lock manager is separate from the physical execution gate. It
retains transaction-level ownership, detects waiting relationships, and
enforces lock waits. A statement can therefore wait even after receiving a
worker. Avoid leaving a transaction open while asking a human a question,
calling a slow external service, or waiting for another application subsystem.

\sourcecode{src/minisql/server/listener.ml}{Bounded connection workers and execution gate dispatch.}
\sourcecode{src/minisql/server/session.ml}{Statement execution and cooperative lock waiting.}
\sourcecode{src/minisql/transaction/lock_manager.ml}{Logical readers, writer, waits, and deadlock handling.}

\section{Deadlines, cancellation, and retry policies}
\index{timeouts}\index{cancellation}\index{retry}

The configured query timeout is an absolute cooperative deadline beginning
at request admission. It covers parsing, gate waiting, execution, and
streamed output. Scans and DML poll at defined boundaries rather than being
preempted at any machine instruction. A client-side timeout and the server's
query deadline are different controls; set them coherently so the application
can receive a useful server diagnostic.

\Needspace{19\baselineskip}
\begin{longtable}{p{.15\linewidth}p{.78\linewidth}}
\toprule
Error & Operational interpretation \\
\midrule
\endhead
9032 & Established logical lock-wait failure. The wait consumes the request
deadline and is counted in timeout statistics. \\
9035 & Administrative cancellation requested for the session's work. \\
9036 & Request deadline reached. \\
9037 & A managed-memory, spill, result-row, or result-byte resource limit was
reached. Increasing a timeout alone cannot fix this limit. \\
9038 & A fenced writer no longer has a valid leadership lease/epoch. Do not
retry writes directly against that retired backend. \\
\bottomrule
\end{longtable}

Python \code{Connection.cancel()} and JDBC \code{Statement.cancel()} use a
separate control connection so cancellation need not compete with the query
already occupying its original socket. Authenticated cancellation requires
database \code{ADMIN}. The administrative chapter explains how to identify
the target session.

For a transient conflict, roll back the transaction, wait a bounded interval,
and retry the whole business operation if it is safe to do so. Limit retries
and preserve the original error for diagnosis. A permanently invalid value,
missing permission, unsupported SQL feature, or lost fencing lease is not a
reason for an unbounded retry loop.

\section{Durability, WAL, and crash recovery}

Changed pages are recorded in a write-ahead log. A successful commit appends
the transaction's records and flushes the WAL before the transaction is
acknowledged as committed. Recovery uses transaction status and page images
to restore committed work and avoid exposing incomplete work. Checkpoints
bound the amount of recovery work by publishing a durable boundary.

Do not delete \path{wal/wal.log}, its durable marker, checkpoint metadata, or
epoch records as a space-saving operation. Their relationship to table pages
is part of recovery correctness. Use supported maintenance and archive
rotation procedures. A crash test that passes on one host is not a guarantee
that every storage device, filesystem, virtual disk, or power-loss scenario
has the same persistence behavior.

MiniSQL exposes only \code{full} durability in its configuration. There is
no documented relaxed-commit setting to trade acknowledged durability for
throughput. If durable writes are slow, first measure the workload and its
transaction size. Grouping related inserts in one transaction can reduce
per-commit overhead, but transactions should remain bounded to avoid large
change sets and long-held locks.

\sourcecode{src/minisql/transaction/transaction.ml}{commit and failure rewind behavior.}
\sourcecode{src/minisql/transaction/recovery.ml}{Recovery validation and replay.}
\sourcecode{src/minisql/transaction/checkpoint.ml}{Durable checkpoint records.}

\section{Practice: verify a transaction, rather than merely typing one}

Create a small disposable table with a primary key and a non-negative quantity
check. Insert a baseline row and record its value. In an explicit transaction,
change the value and read it through the same session; roll back and verify
the original value. Repeat with commit, restart the server, and verify the
new value. Finally create a savepoint, trigger a duplicate-key error, recover
to the savepoint, and commit a valid change.

For each step, predict which value should be visible before executing it.
This exercise connects SQL syntax to session state, rollback semantics, and
durability. Extend it with two sessions only after the single-session result
is understood; otherwise a locking delay can obscure a simpler transaction
mistake.
